\section{A direct PC simulation of the proof DAG}\label{sec:clause-simulation}

Represent a linear clause by its affine true-indicators:
\[
 C=\bigvee_i(g_{C,i}=1),\qquad
 Z(C)=\{b:g_{C,i}(b)=0\text{ for all }i\}.
\]
An equation \(a(b)=\alpha\) has true-indicator \(1+a(b)+\alpha\).
The set \(Z(C)\) is an affine subspace or empty.

Assign to each nonempty clause a fresh complete accuracy-\(h\) block,
and write its product as \(P_C\). Set \(P_\varnothing=1\), using no
block for the empty clause. Fix the whole block family before following
the proof DAG. Clause values \(P_C=0\) will be derived; they are not
included among its axioms.

\begin{lemma}[Prefix identity]\label{lem:prefix}
For every clause block there are polynomials \(U_{C,i}\) satisfying
\begin{equation}\label{eq:prefix}
 1-P_C=\sum_iU_{C,i}g_{C,i},\qquad
 \deg P_C\le2h,\qquad\deg U_{C,i}\le2h-1.
\end{equation}
Its companions have degree at most \(2h+1\).
\end{lemma}
\begin{proof}
Put \(L_u=1-\sum_i r_{ui}g_{C,i}\).
Telescoping the product gives
\[
 1-\prod_{u=1}^hL_u
  =\sum_{u=1}^h(1-L_u)\prod_{w<u}L_w.
\]
Thus one may take
\[
 U_{C,i}=\sum_{u=1}^h r_{ui}\prod_{w<u}L_w.
\]
Each \(L_w\) has degree at most two, giving all the stated bounds.
\end{proof}

\begin{lemma}[Complementary-parity resolution]\label{lem:clause-resolution}
Suppose \(A=(\bigvee_i(a_i=1))\vee(u=1)\) and
\(B=(\bigvee_j(b_j=1))\vee(1-u=1)\), with \(u\) affine.
Let \(T\) have the union of the context inputs \(a_i,b_j\).
If \(P_A,P_B\) have PC derivations through degree \(K\), then \(P_T\)
has a derivation in the same fixed block system through
\(\max\{K,4h+1\}\).
\end{lemma}
\begin{proof}
The context companions \(a_iP_T,b_jP_T\) are axioms of the conclusion
block; a repeated input may use the same axiom.
The prefix identities yield
\begin{align}
 R_A&=P_T-U_{A,u}uP_T
       =P_AP_T+\sum_iU_{A,i}(a_iP_T),\label{eq:resolution-A}\\
 R_B&=P_T-U_{B,1-u}(1-u)P_T
       =P_BP_T+\sum_jU_{B,j}(b_jP_T).\label{eq:resolution-B}
\end{align}
Lemma~\ref{lem:pc-product} derives these through \(\max\{K,4h\}\).
In particular it multiplies the completed lines \(P_A,P_B\), not their
entire earlier derivations.

For an affine \(u=c+\sum_t c_tb_t\), the identity
\[
 u^2-u=\sum_t c_t(b_t^2-b_t)
\]
has a degree-two proof. Therefore
\begin{equation}\label{eq:pivot-correction}
 (1-u)R_A-U_{A,u}P_T(u^2-u)=(1-u)P_T
\end{equation}
is derivable through \(\max\{K,4h+1\}\).
The right side is a completed line of degree at most \(2h+1\).
Multiply it by \(U_{B,1-u}\), costing at most \(4h\), and add
\eqref{eq:resolution-B}. The result is \(P_T\).
All multiplications expand into primitive PC steps by
Lemma~\ref{lem:pc-product}.
\end{proof}

\begin{lemma}[Semantic weakening and tautologies]\label{lem:clause-weakening}
If \(C\models D\) and \(P_C\) has a degree-\(K\) derivation, then
\(P_D\) is derivable through \(\max\{K,4h\}\).
A tautological \(D\) has a direct value derivation through \(2h+1\).
\end{lemma}
\begin{proof}
If \(Z(D)\) is nonempty, then \(Z(D)\subseteq Z(C)\).
Each \(g_{C,i}\) vanishes on \(Z(D)\), so affine linear algebra gives
constants \(\lambda_{ij}\) with
\[
 g_{C,i}=\sum_j\lambda_{ij}g_{D,j}.
\]
To see this directly, choose affine coordinates beginning with a basis
of the equations defining \(Z(D)\). An affine polynomial vanishing when
these coordinates are zero has only those coordinate terms.
Consequently \(g_{C,i}P_D\) is a constant linear combination of
conclusion companions. Now use
\[
 P_D=P_CP_D+\sum_iU_{C,i}(g_{C,i}P_D)
\]
and Lemma~\ref{lem:pc-product}, with ceiling \(\max\{K,4h\}\).

If \(Z(D)\) is empty, its affine equations are inconsistent, and their
span contains one. The same constant combination of their companions
gives \(P_D\), through \(2h+1\).
If the conclusion is empty, its value is one; the preceding formulas
still apply with no conclusion block. In the nonempty-zero-set case,
an implication to the empty clause forces all premise indicators to
be identically zero, so the already derived premise value is one.
\end{proof}

\begin{lemma}[Two-premise semantic inference]\label{lem:affine-cover}
Every sound inference \(A\wedge B\models D\) between linear clauses
over \(\F\) reduces to a tautology case, one semantic weakening, or
two semantic weakenings followed by one complementary-parity
resolution. The last case introduces at most two auxiliary clauses.
\end{lemma}
\begin{proof}
Put \(W=Z(D)\). Soundness says
\(W\subseteq Z(A)\cup Z(B)\).
An empty \(W\) gives a tautology, and containment in either one of the
two sets gives a one-premise weakening.

Otherwise \(W\cap Z(A)\) and \(W\cap Z(B)\) are proper affine
subspaces covering \(W\). Each has at most half its points.
Both must therefore have exactly half, and they must be disjoint.
They are parallel affine hyperplanes of \(W\): two hyperplanes with
different directions would intersect.
Choose an affine \(u\) that is zero on the first and one on the second,
and extend it from \(W\) to the ambient affine space.
Then
\[
 A\models D\vee(u=1),\qquad
 B\models D\vee(u=0).
\]
The two weakenings and a resolution give \(D\).
This argument uses that the full Boolean assignment space is
\(\F^v\).
\end{proof}

For size accounting, replace every clause's falsifying affine system
by a basis, using at most \(v+1\) equations, or by the single
inconsistent equation in the tautology case. Semantic equivalence
preserves the soundness of all inferences. Apply
Lemma~\ref{lem:affine-cover} afterward.
Thus \(S\) proof nodes can be normalized with at most \(3S\) nodes and
polynomial clause descriptions in \(v\).
Unrelated propositional variables may first be fixed to zero; semantic
implication survives restriction, and a constant pivot gives a
weakening. These normalizations justify using line count as the size
measure in Theorem~\ref{thm:main}.

\begin{theorem}[Complete clause-to-PC transfer]\label{thm:clause-transfer}
A refutation of \(\BPHP_n^{n+1}\) with \(S\) nodes yields, for every
integer \(h\ge1\), a complete one-level old-affine family over
\(\mathcal Q_n\) with at most \(3S+\binom m2\) blocks and an ordinary-PC
refutation through
\begin{equation}\label{eq:simulation-degree}
 D=\max\{2h+\ell,4h+1\}.
\end{equation}
After normalization, the total input and coefficient inventory is
polynomial in \(S,n,h\). There is no proof-height restriction.
\end{theorem}
\begin{proof}
Introduce the compact initial linear clauses
\[
 C_{ii'}=\bigvee_{t=1}^\ell(b_{it}+b_{i't}=1).
\]
Each \(C_{ii'}\) semantically implies each label-specific initial clause
for the same pair in \eqref{eq:bit-php}: both rows encoding the same
label falsifies \(C_{ii'}\).
Prepend at most \(\binom m2\) compact clauses and obtain each needed
CNF clause by one weakening, using its existing initial node.
Normalize the remaining proof as above. The total node count is at
most \(3S+\binom m2\).

Assign fresh blocks to its nonempty clauses.
For a compact initial clause set \(g_t=b_{it}+b_{i't}\) and
\(T_t=\prod_{s<t}(1-g_s)\).
Its value has the ordinary identity
\begin{equation}\label{eq:initial-value}
 P_{C_{ii'}}=E_{ii'}P_{C_{ii'}}
                 +\sum_{t=1}^{\ell}T_t(g_tP_{C_{ii'}}).
\end{equation}
The first term uses the old degree-\(\ell\) collision axiom;
the other terms use the block's companions.
Every multiple has degree at most \(2h+\ell\).
This is an NS, hence PC, derivation at that degree.

Follow the normalized source proof in its acyclic order.
Initial values are obtained from \eqref{eq:initial-value} and
Lemma~\ref{lem:clause-weakening}. Each remaining inference is simulated
by Lemmas~\ref{lem:clause-resolution} and \ref{lem:clause-weakening}.
Their premise values are actual earlier derived lines, so the same
global ceiling \eqref{eq:simulation-degree} persists along the DAG.
The final empty clause has value one.
All blocks have inputs affine in the old \(v\) bits and disjoint fresh
coefficients. Their full companion and coefficient-domain families
are present. No balancing or multiplication of an entire preceding
derivation is used.
\end{proof}
